\section{Query Engine}
\label{sec:query-engine}

Both storage engines expose a common query interface. The query pipeline
consists of parsing, planning, optimization, execution, and
canonicalization.

\subsection{Language}
\label{sec:query:language}

The query language supports five operation types:
\begin{itemize}
    \item \textbf{FACT}: retrieve facts matching a pattern.
    \item \textbf{WALK}: traverse relations from an entity.
    \item \textbf{JOIN}: multi-way join across relations.
    \item \textbf{GLUE}: reconstruct global sections from local sections.
\end{itemize}

\subsection{Planning and Optimization}
\label{sec:query:planner}

The logical plan is a tree of operators (Scan, Filter, Join, Project,
Aggregate, Sort, Limit). The optimizer applies:
\begin{itemize}
    \item Predicate pushdown: evaluate filters as early as possible.
    \item Join reordering: order joins by estimated selectivity.
    \item Index selection: choose the most selective index for each scan.
    \item Dead operator elimination: remove unused projections.
\end{itemize}

\subsection{Physical Execution}
\label{sec:query:execution}

Each engine translates the optimized logical plan into engine-specific
physical operators. The KG engine uses index seeks and nested-loop joins.
The Sheaf engine uses context lookups and restriction map traversals.
